\section*{Part A}
\label{sec:part_a}

I think the theme of \textbf{Edgar Allen Poe's} story, \textit{The Premature
Burial's} main theme is how extreme terror can affect one's state of mind.
The narrator's terror is the result of him being obsessed with his fear of being
buried alive.

\begin{center}
  \textit{"I entered into a series of elaborate precautions. Among other things,
  I had the family vault so remodeled as to admit of being readily opened from
  within. The slightest pressure upon a long lever that extended far into the tomb
  would cause the iron portal to fly back. There were arrangements also for the
  free admission of air and light, and convenient receptacles for food and water,
  within immediate reach of the coffin intended for my reception."}
  
  \textit{This coffin was warmly and softly padded, and was provided with a lid,
  fashioned upon the principle of the vault-door, with the addition of springs
  so contrived that the feeblest movement of the body would be sufficient to set
  it at liberty. Besides all this, there was suspended from the roof of the tomb,
  a large bell, the rope of which, it was designed, should extend through a hole
  in the coffin, and so be fastened to one of the hands of the corpse."}
\end{center}

\noindent shows that he is obsessed with this fear and he tries everything he
can to prevent this from happening. The following text uses a sentence structure
to create suspense.

\begin{center}
  \textit{"In all this I endured, there was no physical suffering but of moral
  distress, an infinite suffering. My fancy turned to all things charnel, I talked
  "of worms, of tombs, and epitaphs." I was lost in reveries of death, and the
  idea of premature burial held continual possession of my brain."}

  \textit{"The ghastly Danger to which I was subjected haunted me day and night.
  When night covered the Earth, then, with every horror of thought, I shook—shook
  as the quivering plumes upon the hearse. When Nature could endure wakefulness no
  longer, it was with a struggle that I consented to sleep—for I shuddered to
  reflect that, upon awaking, I might find myself the tenant of a grave."}
\end{center}

It also has some rhythm to it, which adds onto the suspense.

% section part_a (end)

\newpage

\section*{Part B}
\label{sec:part_b}

When I was younger (like 4-5), I was terrified of earth quakes and sink holes. I
would always do research by reading books, websites, etc. I was so obsessed, it
was hard for me to go to sleep because I thought my house will fall down because
of a sink hole appearing. Just like how the character from the story did, I
would always create a plan for each of the scenarios. Just like how the
character's fear took his mind, I eventually started to get more and more
nervous/crazy. But, eventually, I grew out of these two fears.

% section part_b (end)
